\chapter{Postnasal Drip}

\section*{Definition}
Postnasal drip (upper airway cough syndrome) is the sensation of nasal secretions draining posteriorly into the nasopharynx and oropharynx, producing throat clearing, cough, and a sensation of a lump in the throat.

\section*{Pathophysiology}
Excessive or altered mucus production from the nasal mucosa, sinuses, or nasopharynx overwhelms normal mucociliary clearance, causing secretions to accumulate in the pharynx. The sensation may also result from heightened pharyngeal sensitivity rather than actual increased secretion volume.

\section*{Causes}
\begin{itemize}
  \item \textbf{Allergic rhinitis:} Seasonal or perennial allergen triggers causing increased mucus production
  \item \textbf{Non-allergic rhinitis:} Vasomotor rhinitis, gustatory rhinitis, medication-induced (ACE inhibitors, NSAIDs), hormonal (pregnancy, hypothyroidism)
  \item \textbf{Rhinosinusitis:} Acute or chronic sinusitis with mucopurulent drainage
  \item \textbf{Laryngopharyngeal reflux (LPR):} Gastric contents refluxing into the pharynx producing the drip sensation
  \item \textbf{Environmental irritants:} Smoke, dry air, air conditioning, chemical fumes
  \item \textbf{Structural:} Nasal septal deviation, turbinate hypertrophy, nasal polyps
\end{itemize}

\section*{When to See a Doctor}
See your doctor if:
\begin{itemize}
  \item Chronic throat clearing, cough, or globus sensation lasting > 8 weeks
  \item Halitosis (bad breath) or sour taste in the mouth
  \item Hoarseness, dysphagia, or heartburn suggesting LPR
  \item Fever, facial pain, or purulent nasal discharge suggesting sinusitis
  \item Hemoptysis (blood in sputum) or unexplained weight loss
  \item Failure of empiric therapy with antihistamines or nasal steroids
\end{itemize}

\section*{Related Symptoms}
\begin{itemize}
  \item Chronic cough (worse when supine)
  \item Throat clearing and globus sensation (lump in throat)
  \item Sore throat and hoarseness
  \item Halitosis and bad taste
  \item Nausea (from swallowed secretions)
\end{itemize}
\textbf{Important:} Postnasal drip is the most common cause of chronic cough in nonsmokers, yet many people with objective mucus hypersecretion do not complain of drip, suggesting sensitization of pharyngeal mechanoreceptors plays a role.

\section*{Self-Care / Home Management}
\begin{itemize}
  \item Saline nasal irrigation (Neti pot or squeeze bottle) once or twice daily
  \item Increased oral hydration (6--8 glasses of water daily) to thin secretions
  \item Humidifier in the bedroom, especially in dry climates or heated rooms
  \item Elevate head of bed to reduce nocturnal reflux and drip
  \item Avoidance of irritants: smoke, strong perfumes, chemical cleaners
  \item Dietary modifications: limit caffeine, alcohol, chocolate, citrus, and spicy foods (if LPR suspected)
\end{itemize}

\section*{How Your Doctor Will Evaluate You}
\begin{itemize}
  \item Nasal endoscopy to evaluate for purulence, polyps, structural abnormalities, and mucosal inflammation
  \item Allergy testing (skin prick or specific IgE) if allergic rhinitis suspected
  \item CT sinuses without contrast for suspected chronic rhinosinusitis
  \item Laryngoscopy to assess for laryngeal edema, erythema, or cobblestoning suggesting LPR
  \item Trial of empiric therapy (oral antihistamine + intranasal corticosteroid for 4 weeks)
  \item Esophageal pH monitoring or MII-pH for refractory LPR
\end{itemize}

\section*{Treatment Options}
\begin{itemize}
  \item \textbf{Allergic rhinitis:} Intranasal corticosteroids, oral antihistamines, allergen immunotherapy
  \item \textbf{Non-allergic rhinitis:} Intranasal azelastine, ipratropium bromide (for profuse rhinorrhea), or capsaicin
  \item \textbf{Rhinosinusitis:} Antibiotics (if bacterial), saline irrigation, intranasal corticosteroid
  \item \textbf{LPR:} Proton pump inhibitor (omeprazole 20 mg BID), dietary modification, head-of-bed elevation, alginate preparations
  \item \textbf{Mucus thinning agents:} Guaifenesin (600--1200 mg BID) to reduce mucus viscosity
  \item \textbf{Surgery:} Septoplasty, turbinate reduction, endoscopic sinus surgery for structural/CRS-related causes
\end{itemize}

\clearpage
